\documentclass[UTF8,fontset=fandol,zihao=-4]{ctexart}
\usepackage[a4paper,margin=22mm]{geometry}
\usepackage{amsmath,amssymb,booktabs,enumitem}
\usepackage[hidelinks]{hyperref}
\setlength{\parindent}{0pt}
\setlength{\parskip}{0.35em}
\setlist{nosep}
\emergencystretch=2em
\pagestyle{plain}
\ctexset{section={format=\large\bfseries\raggedright,beforeskip=1.2ex,afterskip=.6ex}}
\author{原创微型课程 · 格式示例}
\date{}
\title{模拟卷答案与评分参考}
\begin{document}
\begin{center}{\LARGE\bfseries 模拟卷答案与评分参考\par}\vspace{0.6em}原创微型课程 · 格式示例\end{center}
\vspace{0.5em}
本文件与 20 分模拟卷对应。等价正确方法可按同等步骤给分。
\section*{第 1 题（10 分）}
(1) $\det(A_t)=t-4$，故 $t\ne4$ 时可逆。（行列式 1 分，条件 1 分）

(2) 第二式减去第一式的两倍，得 $(t-4)x_2=0$。当 $t\ne4$ 时，$x_2=0$，$x_1=3$，即 $x=(3,0)^{\mathsf T}$。（消元 1 分，两分量各 1 分）

(3) $t=4$ 时两方程等价，全部解为
\[
x=(3-2s,s)^{\mathsf T},\quad s\in\mathbb R.
\]
（识别依赖 1 分，参数解 2 分）

(4) 无解。两式分别要求 $x_1+2x_2=3$ 与 $2x_1+4x_2=7$，第一式乘二的右端为 6，与 7 矛盾。（结论 1 分，理由 1 分）

\textbf{易错：}行列式为零只说明不可逆，不能直接断言无解。\\
\textbf{核验：}唯一解代入得 $(3,6)^{\mathsf T}$；参数解两式恒成立。\\
\textbf{依据：}S001 第 1--2 节。题目是原创的参数分类练习。

\section*{第 2 题（10 分）}
(1) 特征方程为 $(2-\lambda)^2-1=0$。特征值为 3 和 1，可取
$u_1=(1,1)^{\mathsf T}$、$u_2=(1,-1)^{\mathsf T}$。（方程 1 分，两组特征对各 1 分）

(2) 令 $v=\alpha u_1+\beta u_2$，则 $\alpha+\beta=4$、
$\alpha-\beta=2$，得 $\alpha=3,\beta=1$。（方程组 1 分，系数各 1 分）

(3) 因而
\[
B^n v=3^{n+1}(1,1)^{\mathsf T}+(1,-1)^{\mathsf T}
=\begin{pmatrix}3^{n+1}+1\\3^{n+1}-1\end{pmatrix}.
\]
（利用特征向量作用 2 分，表达式 1 分，检验 1 分）
$n=1$ 得 $(10,8)^{\mathsf T}$，直接相乘 $Bv$ 也得到此结果；
$n=0$ 恢复 $(4,2)^{\mathsf T}$。

\textbf{易错：}系数 3 与特征值的幂要同时保留；特征向量不能取零。\\
\textbf{依据：}S001 第 3 节。S002、S003 的第 2 题仅提供题型参照。\\
\textbf{分值核验：}$2+3+3+2=10$，$3+3+4=10$，全卷 20 分。
\end{document}
